\section{Posterior computation}\label{sec:computation}

Integrating out $\gamma$ gives a posterior density for the remaining
coefficients $\boldsymbol g$. We sample from this density and then draw
$\gamma$ from its conditional truncated normal distribution to recover
the full coefficient vector. We first form the ordinary
Gaussian posterior in \cref{eq:gaussian-posterior} and apply the
appropriate invertible map from
\cref{cor:positive,cor:monotone,cor:convex}, or the common map in
\cref{prop:joint}, to its mean and covariance,
giving the Gaussian laws
\[
 \boldsymbol g\mid\boldsymbol y\sim N(\boldsymbol m_g,\boldsymbol V_g),
 \qquad \gamma\mid\boldsymbol g,\boldsymbol y\sim N(\mu(\boldsymbol g),s^2),
\]
where $\mu$ is affine and $s>0$ is constant in $\boldsymbol g$.
These Gaussian parameters can therefore be prepared before marginal
sampling, together with an invertible factor
$\boldsymbol B\boldsymbol B^T=\boldsymbol V_g$ defining the whitened
coordinates $\boldsymbol g=\boldsymbol m_g+\boldsymbol B\boldsymbol z$,
in which $\boldsymbol z$ has the standard normal law on $\R^m$ under
the ordinary posterior. Mapping through $\boldsymbol B$ preserves the
Gaussian dependence among the remaining coefficients.

At each proposed $\boldsymbol z$, evaluating the marginal density
requires forming $\boldsymbol g$ and its continuum threshold
$\tau_r(\boldsymbol g)$. For the individual cases, this means comparing
both endpoints and all interior stationary candidates from
\cref{cor:positive,cor:monotone,cor:convex}, or obtaining bounds over
the entire interval. The Gaussian log-survival term in
\cref{eq:whitened-target} then weights the standard normal density by
the conditional probability that $\gamma$ satisfies the sign
restriction. Since the observation likelihood is already included
in the ordinary posterior and the overall conditioning probability
is constant in $\boldsymbol z$, neither requires a separate evaluation
during marginal sampling. The monotonicity and
convexity implementations sample this unrestricted $m$-dimensional
density by Hamiltonian Monte Carlo, with the numerical threshold
evaluation described below.

For each retained $\boldsymbol g$, we recover the separated coefficient
by drawing $\gamma$ from
$N(\mu(\boldsymbol g),s^2)$ truncated to
$[\tau_r(\boldsymbol g),\infty)$, using an independent
$U\sim\operatorname{Uniform}(0,1)$
and the inverse-survival formula in \cref{eq:reconstruction}, which
retains its conditional uncertainty. Applying the corresponding inverse
map in \cref{tab:coefficient-reconstruction} then recovers the original
coefficients of $f_{\boldsymbol a}$.
With the exact threshold and draws from the
exact marginal law, these two stages recover the full conditioned
coefficient law, so evaluating the expansion and its derivatives gives
joint posterior curves from which to compute means and pointwise
credible intervals.

For a signed joint restriction satisfying \cref{prop:joint}, the common
map determines both the Gaussian preparation and every quotient
$\eta_r(x;\boldsymbol g)$. At each proposed $\boldsymbol g$, we
evaluate the required continuum and endpoint bounds and take their
maximum to obtain $\tau_A^\varepsilon(\boldsymbol g)$. Since these
bounds all concern the same $\gamma$, a single survival probability
$\Phi^{\mathrm c}((\tau_A^\varepsilon(\boldsymbol g)-\mu(\boldsymbol g))/s)$
weights the marginal density for the entire intersection.
Reconstruction uses this joint lower bound in the truncated normal
and then inverts the common coefficient map. The Gaussian parameters
must be computed under that map, and the individual thresholds in
\cref{tab:cases}, which use different maps, cannot be combined
directly. The supplementary material specifies the HSGP map and
the continuum and endpoint bounds for each joint combination.

In numerical evaluation, a successful threshold calculation returns an
enclosure whose upper endpoint $\widetilde\tau_r$ satisfies
\[
 \tau_r(\boldsymbol g)\leq\widetilde\tau_r(\boldsymbol g)
       \leq\tau_r(\boldsymbol g)+\delta
\]
for the supplied coefficients and a chosen tolerance $\delta>0$.
If marginal weighting and conditional reconstruction use the same
numerical threshold and the enclosure bound holds almost everywhere
under the ordinary posterior, \cref{prop:error} controls the total
variation distance between the resulting target distribution and the
posterior conditioned on the exact sign event.
The conditional standard deviation $s$ provides a natural scale for
choosing $\delta$, because $\delta/s$ bounds the error in the
standardized truncation point. For a prescribed total variation
error $\rho\in(0,1)$, \cref{prop:error} gives the sufficient choice
\[
 \delta\leq\rho\,s\sqrt{2\pi}\,P_{\boldsymbol y}(\mathcal C_r).
\]
In practice, a valid positive lower bound on the conditioning
probability can replace $P_{\boldsymbol y}(\mathcal C_r)$ on the
right-hand side. Without such a bound, a small value of $\delta/s$
controls the standardized threshold error but does not by itself
certify a small posterior error, particularly when the shape event
is rare.

Stable evaluation uses log-survival and inverse-survival calculations
to avoid subtracting nearly equal probabilities in the tails. For
gradient-based sampling, the derivative implementations evaluate
approximate envelope gradients at numerically selected maximizing
locations. At a unique maximum these gradients approximate derivatives of the
exact threshold, while the sampled log values use the rounded upper
threshold. The supplementary material explains how the numerical
values and approximate gradients enter the Stan No-U-Turn sampler
(NUTS), alongside the compact algorithm, Gaussian block formulas,
threshold enclosures, precision controls, and stable reconstruction.
